\documentclass{zhvt-classic}

% 使用本地源流明体（GenRyuMin2 TW）按文件路径加载，不依赖系统字体名
% 版框（内框／界栏／书口侧栏／鱼尾）改用与版心印章同一印泥棕色：
%   印章图片 yshs.png 实色 = RGB(135,68,52) = HTML 874434
\usepackage{xcolor}
\usepackage{tikz}
% sealbrown 双注册：xcolor 供 TikZ/pgf，L3 供 cls 版框
\definecolor{sealbrown}{HTML}{874434}
\ExplSyntaxOn
\color_set:nnn { sealbrown } { HTML } { 874434 }
\ExplSyntaxOff

\zhvtset{
  grid color       = sealbrown,
  main font        = GenRyuMin2TW-SB,
  main font path   = C:/Users/Administrator/texmf/fonts/opentype/public/GenRyuMin2TW/,
jiazhu font      = GenRyuMin2TW-M,
  jiazhu font path = C:/Users/Administrator/texmf/fonts/opentype/public/GenRyuMin2TW/,
  % 页码框容纳 4 字：正常页「九十七」+前置空 1 格=4 行恰好；
  % 100+ 页「一百零一」共 5 行超出，末字「一/二/三」被框裁掉（2026-10-04 实测），故放宽到 5
  shukou page number chars = 5,
  % 书口书名框高从默认 6 字减到 4 字，把省出的 2 字高拨给 chapter 框，
  % 使 6 字书名+卷名放进 9 字高的 chapter 框后不再压线（2026-10-05）
  shukou title chars = 4,
}

\title{史记}[留侯世家]
\author{汉·司马迁}
\maker{闲叟漫记·重製}


% ===== 索隐述赞样式（2026-11-08 四次定稿 · 愚大人亲手改定参数）=====
% ⚠ 竖排前提：自定义字体家族必须带 RawFeature={vertical:+vert}；
%   simhei.ttf 加 +vhal 会整体加载失败静默回退，只写 +vert；simfang 同理。
\newCJKfontfamily{\shuzanfont}[Path=C:/Windows/Fonts/,RawFeature={vertical:+vert}]{simhei.ttf}
\newCJKfontfamily{\shuzanbody}[Path=C:/Windows/Fonts/,RawFeature={vertical:+vert}]{simfang.ttf}
% 题头：152x30pt 圆角 6pt 棕底白字；baseline (0,2) 使整框落入界行格子不压版框线；
% 白字节点 x 不补偿、y +8.2pt（像素级墨迹实测 dx/dy≈0，L/R 留白约 20pt）。
\newcommand{\shuzanhead}[1]{%
  \begin{tikzpicture}[baseline={(0,2)}]
    \fill[sealbrown, rounded corners=6pt]
      (-76pt,-15pt) rectangle (76pt,15pt);
    \node[text=white, inner sep=0pt, anchor=center] at (0pt,8.2pt)
      {\shuzanfont\fontsize{23}{23}\selectfont #1};
  \end{tikzpicture}%
}
% ★ 修正书口卷名位置：chapter 框起始位置改回 cls 默认值，避免全局影响。
%   只在每个卷首页通过 \juan 命令里的空命令微调该页位置（安全，不会跨页污染）。
\ExplSyntaxOn
\AtBeginDocument{
  \dim_set:Nn \l__zhvt_shukou_chapter_title_start_dim
    { \l__zhvt_shukou_title_height_dim + \l__zhvt_font_size_dim * 2 }
  \dim_set:Nn \l__zhvt_shukou_page_number_start_dim
    { \l__zhvt_shukou_chapter_title_start_dim + \l__zhvt_shukou_chapter_title_height_dim }
}
\ExplSyntaxOff

\begin{document}

% ★★ 版心（书口）排布覆写：把书名放到鱼尾之下（2026-10-04）
%
%   cls 默认版心自上而下：书名（title 静态框，在鱼尾之上）→ 鱼尾 → 篇名（chapter 动态框）→ 页码
%   本模板改为：          鱼尾（位置不动）→ 书名 → 篇名 → 页码分割线 → 页码
%
%   手法：把 cls 放在 title 静态框里的书名，改放进 chapter 动态框（鱼尾之下）；
%         title 静态框留空。书名取 \l__zhvt_title_str，即 \title{...} 的第一参数，
%         所以换书时只改 \title，不必在此重复写书名。
%
%   ⚠ 三条必须遵守的写法（2026-10-04 实测）：
%     ① 汉字必须用 \hbox{...} 包住。flowfram 会把框内容装进一个宽 14.4mm（= 一个 line_sep
%        = 书口栏宽）的纵向 minipage 并按该宽度折行，而整页靠 PDF /Rotate 90 显示、
%        框内文字本身不变向 —— 不包 \hbox 时「围炉夜话」4 字会被逐字折成 4 行，
%        在版面上就成了横排 4 字并横向溢出版框。
%     ② 竖排下**一个 \hbox 就是一整列**。书名与篇名要同处一列，必须放进同一个 \hbox；
%        分成两个 \hbox 时，后一个会被当成新的一列、甩到隔壁栏去
%        （实测「序」跑到书口栏外侧 14.4mm 的正文列上）。中间的 \quad 在框内是横向空白，
%        显示出来正好是书名与篇名之间的纵向间隔。
%     ③ ★ 内容盒要写成 **与框等宽的定宽盒** `\hbox to \l__zhvt_shukou_chapter_title_height_dim{…\hfil}`。
%        实测：chapter 框是把内容盒沿纵向**居中**摆放的，内容宽度一变，书名就跟着漂：
%        自然宽度的 \hbox（书名 38.4mm）落在 x=157.2mm，加篇名后（57.6mm）落在 x=152.4mm，
%        目录页（无篇名）与正文页的书名就差半格。做成等宽盒后，无论有没有篇名，
%        内容固定从框顶起排，实测两页均落在 x=142.8mm，不再漂。
\ExplSyntaxOn
% 把书口 chapter 框（书名+篇名）起始位置下移 1 字高，避免 4 字卷名顶到鱼尾
% ⚠ cls 的 \__zhvt_calculate_dimensions: 在 \AtBeginDocument 里执行，会重算该 dim；
%   必须把覆写挂到 begin document 之后，否则会被二次覆盖（2026-10-05 实测）。
% (旧方案已移到导言区，直接修改 chapter 框起始位置)
% 卷标题命令 #1=书口篇名(鱼尾) #2=正文卷名（含卷序，卷名内空格用 ~）
%
% ★★ 目录页两项修订（2026-10-05，愚大人指出）★★
%   ① 页码起点统一对齐：
%      原状 —— 目录项写成 \par\hspace{4em}卷名\hspace{4em}页码，页码起点 = 卷名实宽 + 4em。
%      卷名实宽随字数变：「卷一~梁惠王上」6 字 +~、「卷七~离娄上」5 字 +~，
%      差 1 个字位 → 卷七~卷十 4 条的页码整体比其余 10 条高 27.21pt（一个字），看着上下不齐。
%      改法 —— 卷名补全角空格（宽 = 1 字 = 27.215pt）补足到 **7 字符位**再写进目录，
%      14 条页码起点遂全部落在同一条竖线上（起点对齐；末字随页码字数自然参差）。
%      ⚠ 只对目录里的副本 \l_tmpc_tl 补，正文卷名仍用 \l_tmpa_tl 原样 —— 正文像素不受影响。
%   ② 卷名整体上移半字位：前导 \hspace 由 4em 改 3.5em。
%      实测原状目录块在版框内偏下（天头留白 139.6pt / 地脚留白 110.7pt），
%      上移 0.5em = 13.6pt 后上下基本均分，视觉居中。
\NewDocumentCommand { \juan } { mm } {
	\clearpage
	\tl_set:Nn \l__zhvt_chapter_title_tl { #1 }
	\tl_set:Nn \l_tmpa_tl { #2 }
	\tl_set_eq:NN \l_tmpc_tl \l_tmpa_tl
	\int_step_inline:nn { 7 - \str_count:n { #2 } }
	  { \tl_put_right:Nn \l_tmpc_tl { 　 } }
	\tl_set:Nx \l_tmpb_tl { \zhnum { page } }
	\tl_put_left:Nn \l_tmpb_tl { \hspace { 4em } }
	\tl_put_left:Nn \l_tmpb_tl { \l_tmpc_tl }
	\tl_put_left:Nn \l_tmpb_tl { \par \hspace { 3.5em } }
	\addtocontents { toc } { \l_tmpb_tl }
	\skip_horizontal:n { \l__zhvt_font_size_dim * 4 }
	\l_tmpa_tl
	\par
	\thispagestyle { empty }
}
\cs_set_protected:Npn \__zhvt_set_shukou_content: {
	% ★ 鱼尾上方留白（title 静态框）嵌入闲叟钤印：
	%   整页 /Rotate 90，故此处把图用 \rotatebox{90} 转正为竖章；
	%   尺寸卡在版心条幅内（书内约 7mm 宽 × 33.8mm 高，留余量不压鱼尾）。
	\setstaticcontents * { title } {
		\vbox_to_ht:nn { \l__zhvt_baseline_skip_dim } {
			\vskip 14.2pt
			\hbox_to_wd:nn { \l__zhvt_shukou_title_height_dim } {
				\hfil
				\rotatebox [ origin = c ] { 90 } { \includegraphics [ width = 7mm , height = 33.8mm ] { yshs.png } }
				\hfil
			}
		}
	}
% ★ 版心（书口）统一改法（2026-11-16 愚大人定）：18pt＋框内居中＋书口文字同版框色
%   ① 原 27.2pt 时 7 字长篇名（如「廉颇蔺相如列传」）末字溢出 chapter 框、压住页码；
%      改 \fontsize{18pt}{18pt}，前后 \hfil 令字串在框内居中（不随篇名字数漂移）。
%   ② \vbox{\vskip 8.2pt} 定量下移，抵消缩字号后栏宽方向 8.2pt 上偏（像素实测 y 中心 543.3 = 原版）。
%      ⚠ \raisebox / \vspace* 在 flowfram 框内无效；换字号须同步换补偿：16pt→10.1、19pt→7.4。
%   ③ chapter 与 page~number 均置 \color{sealbrown}（HTML 874434），与版框／界行同色；正文仍黑。
%   ⚠ 页码字位、钤印、鱼尾、界行、正文一律未动；五部合订本同此口径。
	\setdynamiccontents * { chapter }
	{ \color{sealbrown} \vbox { \vskip 8.2pt \hbox to \l__zhvt_shukou_chapter_title_height_dim
		{ \hfil \fontsize{18pt}{18pt}\selectfont \l__zhvt_title_str \quad \l__zhvt_chapter_title_tl \hfil } } }
	\setdynamiccontents * { page~number } { \color{sealbrown} \quad \zhnum { page } }
}
\__zhvt_set_shukou_content:
\ExplSyntaxOff

\maketitle

% 单篇，无多卷目录，不渲染目录页

\juan{留侯世家}{卷五十五~留侯世家第二十五}
\hspace*{2em}留侯张良者。其先韩人也。大父开地。相韩昭侯。宣惠王。襄哀王。父平。相釐王。悼惠王。悼惠王二十三年。平卒。卒二十岁。秦灭韩。良年少。未宦事韩。韩破。良家僮三百人。弟死不葬。悉以家财求客刺秦王。为韩报仇。以大父。父五世相韩故。\\
\hspace*{2em}良尝学礼淮阳。东见仓海君。得力士。为铁椎重百二十斤。秦皇帝东游。良与客狙击秦皇帝博浪沙中。误中副车。秦皇帝大怒。大索天下。求贼甚急。为张良故也。良乃更名姓。亡匿下邳。\\
\hspace*{2em}良尝闲从容步游下邳圯上。有一老父。衣褐。至良所。直堕其履圯下。顾谓良曰。﹁孺子。下取履。﹂良鄂然。欲殴之。为其老。彊忍。下取履。父曰。﹁履我。﹂良业为取履。因长跪履之。父以足受。笑而去。良殊大惊。随目之。父去里所。复还。曰。﹁孺子可教矣。后五日平明。与我会此。﹂良因怪之。跪曰。﹁诺。﹂五日平明。良往。父已先在。怒曰。﹁与老人期。后。何也。﹂去。曰。﹁后五日早会。﹂五日鸡鸣。良往。父又先在。复怒曰。﹁后。何也。﹂去。曰。﹁后五日复早来。﹂五日。良夜未半往。有顷。父亦来。喜曰。﹁当如是。﹂出一编书。曰。﹁读此则为王者师矣。后十年兴。十三年孺子见我济北。穀城山下黄石即我矣。﹂遂去。无他言。不复见。旦日视其书。乃太公兵法也。良因异之。常习诵读之。\\
\hspace*{2em}居下邳。为任侠。项伯尝杀人。从良匿。\\
\hspace*{2em}后十年。陈涉等起兵。良亦聚少年百馀人。景驹自立为楚假王。在留。良欲往从之。道遇沛公。沛公将数千人。略地下邳西。遂属焉。沛公拜良为厩将。良数以太公兵法说沛公。沛公善之。常用其策。良为他人言。皆不省。良曰。﹁沛公殆天授。﹂故遂从之。不去见景驹。\\
\hspace*{2em}及沛公之薛。见项梁。项梁立楚怀王。良乃说项梁曰。﹁君已立楚后。而韩诸公子横阳君成贤。可立为王。益树党。﹂项梁使良求韩成。立以为韩王。以良为韩申徒。与韩王将千馀人西略韩地。得数城。秦辄复取之。往来为游兵颍川。\\
\hspace*{2em}沛公之从雒阳南出轘辕。良引兵从沛公。下韩十馀城。击破杨熊军。沛公乃令韩王成留守阳翟。与良俱南。攻下宛。西入武关。沛公欲以兵二万人击秦峣下军。良说曰。﹁秦兵尚彊。未可轻。臣闻其将屠者子。贾竖易动以利。原沛公且留壁。使人先行。为五万人具食。益为张旗帜诸山上。为疑兵。令郦食其持重宝啗秦将。﹂秦将果畔。欲连和俱西袭咸阳。沛公欲听之。良曰。﹁此独其将欲叛耳。恐士卒不从。不从必危。不如因其解击之。﹂沛公乃引兵击秦军。大破之。北至蓝田。再战。秦兵竟败。遂至咸阳。秦王子婴降沛公。\\
\hspace*{2em}沛公入秦宫。宫室帷帐狗马重宝妇女以千数。意欲留居之。樊哙谏沛公出舍。沛公不听。良曰。﹁夫秦为无道。故沛公得至此。夫为天下除残贼。宜缟素为资。今始入秦。即安其乐。此所谓﹃助桀为虐﹄。且﹃忠言逆耳利于行。良药苦口利于病﹄。愿沛公听樊哙言。﹂沛公乃还军霸上。\\
\hspace*{2em}项羽至鸿门下。欲击沛公。项伯乃夜驰入沛公军。私见张良。欲与俱去。良曰。﹁臣为韩王送沛公。今事有急。亡去不义。﹂乃具以语沛公。沛公大惊。曰。﹁为将柰何。﹂良曰。﹁沛公诚欲倍项羽邪。﹂沛公曰。﹁鲰生教我距关无内诸侯。秦地可尽王。故听之。﹂良曰。﹁沛公自度能卻项羽乎。﹂沛公默然良久。曰。﹁固不能也。今为柰何。﹂良乃固要项伯。项伯见沛公。沛公与饮为寿。结宾婚。令项伯具言沛公不敢倍项羽。所以距关者。备他盗也。及见项羽后解。语在项羽事中。\\
\hspace*{2em}汉元年正月。沛公为汉王。王巴蜀。汉王赐良金百溢。珠二斗。良具以献项伯。汉王亦因令良厚遗项伯。使请汉中地。项王乃许之。遂得汉中地。汉王之国。良送至襃中。遣良归韩。良因说汉王曰。﹁王何不烧绝所过栈道。示天下无还心。以固项王意。﹂乃使良还。行。烧绝栈道。\\
\hspace*{2em}良至韩。韩王成以良从汉王故。项王不遣成之国。从与俱东。良说项王曰。﹁汉王烧绝栈道。无还心矣。﹂乃以齐王田荣反。书告项王。项王以此无西忧汉心。而发兵北击齐。\\
\hspace*{2em}项王竟不肯遣韩王。乃以为侯。又杀之彭城。良亡。间行归汉王。汉王亦已还定三秦矣。复以良为成信侯。从东击楚。至彭城。汉败而还。至下邑。汉王下马踞鞍而问曰。﹁吾欲捐关以东等弃之。谁可与共功者。﹂良进曰。﹁九江王黥布。楚枭将。与项王有郄。彭越与齐王田荣反梁地。此两人可急使。而汉王之将独韩信可属大事。当一面。即欲捐之。捐之此三人。则楚可破也。﹂汉王乃遣随何说九江王布。而使人连彭越。及魏王豹反。使韩信将兵击之。因举燕。代。齐。赵。然卒破楚者。此三人力也。\\
\hspace*{2em}张良多病。未尝特将也。常为画策。时时从汉王。\\
\hspace*{2em}汉三年。项羽急围汉王荥阳。汉王恐忧。与郦食其谋桡楚权。食其曰。﹁昔汤伐桀。封其后於杞。武王伐纣。封其后於宋。今秦失德弃义。侵伐诸侯社稷。灭六国之后。使无立锥之地。陛下诚能复立六国后世。毕已受印。此其君臣百姓必皆戴陛下之德。莫不乡风慕义。原为臣妾。德义已行。陛下南乡称霸。楚必敛衽而朝。﹂汉王曰。﹁善。趣刻印。先生因行佩之矣。﹂\\
\hspace*{2em}食其未行。张良从外来谒。汉王方食。曰。﹁子房前。客有为我计桡楚权者。﹂其以郦生语告。曰。﹁於子房何如。﹂良曰。﹁谁为陛下画此计者。陛下事去矣。﹂汉王曰。﹁何哉。﹂张良对曰。﹁臣请藉前箸为大王筹之。﹂曰。﹁昔者汤伐桀而封其后於杞者。度能制桀之死命也。今陛下能制项籍之死命乎。﹂曰。﹁未能也。﹂﹁其不可一也。武王伐纣封其后於宋者。度能得纣之头也。今陛下能得项籍之头乎。﹂曰。﹁未能也。﹂﹁其不可二也。武王入殷。表商容之闾。释箕子之拘。封比干之墓。今陛下能封圣人之墓。表贤者之闾。式智者之门乎。﹂曰。﹁未能也。﹂﹁其不可三也。发钜桥之粟。散鹿台之钱。以赐贫穷。今陛下能散府库以赐贫穷乎。﹂曰。﹁未能也。﹂﹁其不可四矣。殷事已毕。偃革为轩。倒置干戈。覆以虎皮。以示天下不复用兵。今陛下能偃武行文。不复用兵乎。﹂曰。﹁未能也。﹂﹁其不可五矣。休马华山之阳。示以无所为。今陛下能休马无所用乎。﹂曰。﹁未能也。﹂﹁其不可六矣。放牛桃林之阴。以示不复输积。今陛下能放牛不复输积乎。﹂曰。﹁未能也。﹂﹁其不可七矣。且天下游士离其亲戚。弃坟墓。去故旧。从陛下游者。徒欲日夜望咫尺之地。今复六国。立韩。魏。燕。赵。齐。楚之后。天下游士各归事其主。从其亲戚。反其故旧坟墓。陛下与谁取天下乎。其不可八矣。且夫楚唯无彊。六国立者复桡而从之。陛下焉得而臣之。诚用客之谋。陛下事去矣。﹂汉王辍食吐哺。骂曰。﹁竖儒。几败而公事。﹂令趣销印。\\
\hspace*{2em}汉四年。韩信破齐而欲自立为齐王。汉王怒。张良说汉王。汉王使良授齐王信印。语在淮阴事中。\\
\hspace*{2em}其秋。汉王追楚至阳夏南。战不利而壁固陵。诸侯期不至。良说汉王。汉王用其计。诸侯皆至。语在项籍事中。\\
\hspace*{2em}汉六年正月。封功臣。良未尝有战斗功。高帝曰。﹁运筹策帷帐中。决胜千里外。子房功也。自择齐三万户。﹂良曰。﹁始臣起下邳。与上会留。此天以臣授陛下。陛下用臣计。幸而时中。臣原封留足矣。不敢当三万户。﹂乃封张良为留侯。与萧何等俱封。\\
\hspace*{2em}上已封大功臣二十馀人。其馀日夜争功不决。未得行封。上在雒阳南宫。从复道望见诸将往往相与坐沙中语。上曰。﹁此何语。﹂留侯曰。﹁陛下不知乎。此谋反耳。﹂上曰。﹁天下属安定。何故反乎。﹂留侯曰。﹁陛下起布衣。以此属取天下。今陛下为天子。而所封皆萧。曹故人所亲爱。而所诛者皆生平所仇怨。今军吏计功。以天下不足遍封。此属畏陛下不能尽封。恐又见疑平生过失及诛。故即相聚谋反耳。﹂上乃忧曰。﹁为之柰何。﹂留侯曰。﹁上平生所憎。群臣所共知。谁最甚者。﹂上曰。﹁雍齿与我故。数尝窘辱我。我欲杀之。为其功多。故不忍。﹂留侯曰。﹁今急先封雍齿以示群臣。群臣见雍齿封。则人人自坚矣。﹂於是上乃置酒。封雍齿为什方侯。而急趣丞相。御史定功行封。群臣罢酒。皆喜曰。﹁雍齿尚为侯。我属无患矣。﹂\\
\hspace*{2em}刘敬说高帝曰。﹁都关中。﹂上疑之。左右大臣皆山东人。多劝上都雒阳。﹁雒阳东有成皋。西有殽黾。倍河。向伊雒。其固亦足恃。﹂留侯曰。﹁雒阳虽有此固。其中小。不过数百里。田地薄。四面受敌。此非用武之国也。夫关中左殽函。右陇蜀。沃野千里。南有巴蜀之饶。北有胡苑之利。阻三面而守。独以一面东制诸侯。诸侯安定。河渭漕輓天下。西给京师。诸侯有变。顺流而下。足以委输。此所谓金城千里。天府之国也。刘敬说是也。﹂於是高帝即日驾。西都关中。\\
\hspace*{2em}留侯从入关。留侯性多病。即道引不食穀。杜门不出岁馀。\\
\hspace*{2em}上欲废太子。立戚夫人子赵王如意。大臣多谏争。未能得坚决者也。吕后恐。不知所为。人或谓吕后曰。﹁留侯善画计筴。上信用之。﹂吕后乃使建成侯吕泽劫留侯。曰。﹁君常为上谋臣。今上欲易太子。君安得高枕而卧乎。﹂留侯曰。﹁始上数在困急之中。幸用臣筴。今天下安定。以爱欲易太子。骨肉之间。虽臣等百馀人何益。﹂吕泽彊要曰。﹁为我画计。﹂留侯曰。﹁此难以口舌争也。顾上有不能致者。天下有四人。四人者年老矣。皆以为上慢侮人。故逃匿山中。义不为汉臣。然上高此四人。今公诚能无爱金玉璧帛。令太子为书。卑辞安车。因使辩士固请。宜来。来。以为客。时时从入朝。令上见之。则必异而问之。问之。上知此四人贤。则一助也。﹂於是吕后令吕泽使人奉太子书。卑辞厚礼。迎此四人。四人至。客建成侯所。\\
\hspace*{2em}汉十一年。黥布反。上病。欲使太子将。往击之。四人相谓曰。﹁凡来者。将以存太子。太子将兵。事危矣。﹂乃说建成侯曰。﹁太子将兵。有功则位不益太子。无功还。则从此受祸矣。且太子所与俱诸将。皆尝与上定天下枭将也。今使太子将之。此无异使羊将狼也。皆不肯为尽力。其无功必矣。臣闻﹃母爱者子抱﹄。今戚夫人日夜待御。赵王如意常抱居前。上曰﹃终不使不肖子居爱子之上﹄。明乎其代太子位必矣。君何不急请吕后承间为上泣言。﹃黥布。天下猛将也。善用兵。今诸将皆陛下故等夷。乃令太子将此属。无异使羊将狼。莫肯为用。且使布闻之。则鼓行而西耳。上虽病。彊载辎车。卧而护之。诸将不敢不尽力。上虽苦。为妻子自彊。﹄﹂於是吕泽立夜见吕后。吕后承间为上泣涕而言。如四人意。上曰。﹁吾惟竖子固不足遣。而公自行耳。﹂於是上自将兵而东。群臣居守。皆送至灞上。留侯病。自彊起。至曲邮。见上曰。﹁臣宜从。病甚。楚人剽疾。原上无与楚人争锋。﹂因说上曰。﹁令太子为将军。监关中兵。﹂上曰。﹁子房虽病。彊卧而傅太子。﹂是时叔孙通为太傅。留侯行少傅事。\\
\hspace*{2em}汉十二年。上从击破布军归。疾益甚。愈欲易太子。留侯谏。不听。因疾不视事。叔孙太傅称说引古今。以死争太子。上详许之。犹欲易之。及燕。置酒。太子侍。四人从太子。年皆八十有馀。须眉皓白。衣冠甚伟。上怪之。问曰。﹁彼何为者。﹂四人前对。各言名姓。曰东园公。角里先生。绮里季。夏黄公。上乃大惊。曰。﹁吾求公数岁。公辟逃我。今公何自从吾兒游乎。﹂四人皆曰。﹁陛下轻士善骂。臣等义不受辱。故恐而亡匿。窃闻太子为人仁孝。恭敬爱士。天下莫不延颈欲为太子死者。故臣等来耳。﹂上曰。﹁烦公幸卒调护太子。﹂\\
\hspace*{2em}四人为寿已毕。趋去。上目送之。召戚夫人指示四人者曰。﹁我欲易之。彼四人辅之。羽翼已成。难动矣。吕后真而主矣。﹂戚夫人泣。上曰。﹁为我楚舞。吾为若楚歌。﹂歌曰。﹁鸿鹄高飞。一举千里。羽翮已就。横绝四海。横绝四海。当可柰何。虽有矰缴。尚安所施。﹂歌数阕。戚夫人嘘唏流涕。上起去。罢酒。竟不易太子者。留侯本招此四人之力也。\\
\hspace*{2em}留侯从上击代。出奇计马邑下。及立萧何相国。所与上从容言天下事甚众。非天下所以存亡。故不著。留侯乃称曰。﹁家世相韩。及韩灭。不爱万金之资。为韩报雠彊秦。天下振动。今以三寸舌为帝者师。封万户。位列侯。此布衣之极。於良足矣。愿弃人间事。欲从赤松子游耳。﹂乃学辟穀。道引轻身。会高帝崩。吕后德留侯。乃彊食之。曰。﹁人生一世间。如白驹过隙。何至自苦如此乎。﹂留侯不得已。彊听而食。\\
\hspace*{2em}后八年卒。谥为文成侯。子不疑代侯。\\
\hspace*{2em}子房始所见下邳圯上老父与太公书者。后十三年从高帝过济北。果见穀城山下黄石。取而葆祠之。留侯死。并葬黄石。每上冢伏腊。祠黄石。\\
\hspace*{2em}留侯不疑。孝文帝五年坐不敬。国除。\\
\hspace*{2em}太史公曰。学者多言无鬼神。然言有物。至如留侯所见老父予书。亦可怪矣。高祖离困者数矣。而留侯常有功力焉。岂可谓非天乎。上曰。﹁夫运筹筴帷帐之中。决胜千里外。吾不如子房。﹂余以为其人计魁梧奇伟。至见其图。状貌如妇人好女。盖孔子曰。﹁以貌取人。失之子羽。﹂留侯亦云。\\

% ===== 索隐述赞（司马贞《史记索隐》）· 统一样式 =====
% 题头「索隐述赞曰」五字并入反白框；正文每四句一列，定宽盒 \hfil 居列中；
% \fontsize{24}{41} 列距对齐全书界行（121+41k）；\clearpage 使述赞独占一页。
\clearpage
\shuzanhead{索隐述赞曰}\linebreak
{\shuzanbody\fontsize{24}{41}\selectfont
\hbox to \hsize{\hfil 留侯倜傥。志怀愤惋。五代相韩。一朝归汉。\hfil}\linebreak
\hbox to \hsize{\hfil 进履宜假。运筹神算。横阳既立。申徒作扞。\hfil}\linebreak
\hbox to \hsize{\hfil 灞上扶危。固陵静乱。人称三杰。辩推八难。\hfil}\linebreak
\hbox to \hsize{\hfil 赤松原游。白驹难绊。嗟彼雄略。曾非魁岸。\hfil}%
\par}
\end{document}
